\documentclass[12pt,a4paper]{article}
\pagestyle{plain}
\date{}
%-------------------------------------------------------------------------------------------------------------------------------------------------
%%%%%%%%%%%%%%%%%%%%%%%%%%%%%%%%%%%%%%%%%%%%%%%%%%%%%%%%%%%%%%%%%%%%%%%%%%%%
%Packages
%%%%%%%%%%%%%%%%%%%%%%%%%%%%%%%%%%%%%%%%%%%%%%%%%%%%%%%%%%%%%%%%%%%%%%%%%%%%
\usepackage{amsmath,amssymb,amscd,amsthm}

\usepackage{multicol,longtable,tabularx}
%\usepackage{array,epic,eepic}
\usepackage{enumerate}

\usepackage{latexsym}

%\usepackage[scanall]{psfrag}
%\usepackage{PSFigure,epsfig,graphicx}
\usepackage{float}


\usepackage{eufrak}
\usepackage{ifthen}
\usepackage{verbatim}
\usepackage{calc}
%\usepackage{refcheck}

%%%%%%%%%%%%%%%%%%%%%%%%%%%%%%%%%%%%%%%%%%%%%%%%%%%%%%%%%%%%%%%%%%%%%%%%%%%%
% Theoremstyle:
%%%%%%%%%%%%%%%%%%%%%%%%%%%%%%%%%%%%%%%%%%%%%%%%%%%%%%%%%%%%%%%%%%%%%%%%%%%%
\newtheorem{theo}{Theorem}[section]
\newtheorem{prop}[theo]{Proposition}
\newtheorem{lemma}[theo]{Lemma}
\newtheorem{cor}[theo]{Corollary}
%\theoremstyle{definition}
\newtheorem{defi}[theo]{Definition}
%\theoremstyle{definition}
\newtheorem{eg}[theo]{Examples}
\newtheorem{es}[theo]{Example}
\theoremstyle{definition}
\newtheorem{oss}[theo]{Remark}
\newtheorem{remarks}[theo]{Remarks}
%%teoremi e simili----------------------------------------------------------------------------------------------------------------------------------

%\newcommand{\proof}{\noindent{\em Proof. } \,}
\newcommand{\proofend}{\hfill{$\square$}}
%numerazione----------------------------------------------------------------------------------------------------------------------------------------
\numberwithin{equation}{section} \numberwithin{figure}{section}
\numberwithin{table}{section}
%separazione parole----------------------------------------------------------------------------------------------------------------------------------
%\hyphenation{co-mo-no-to-ni-ci-ty}


%%%%%%%%%%%%%%%%%%%%%%%%%%%%%%%%%%%%%%%%%%%%%%%%%%%%%%%%%%%%%%%%%%%%%%%%%%%%%%%%%%%%%%%%%%%%%%%%%%%%%%%%%%%%%%%%%%%%%%%%%%%%%%%%%%%%%%%%%%%%%%%%%%%
\begin{document}
\begin{center}
Esame Analisi Matematica (A-L)\\
Corso di Laurea in Informatica\\
09-04-26
\end{center}

\ \\
1.  \indent Si risolvano i seguenti limiti:
\begin{enumerate}
\item[(a)] $\lim\limits_{x \to -\infty} \sqrt{x^2-4x+1}+x;$
\item[(b)] $\lim\limits_{n \to +\infty} \dfrac{n (e^{2/n^2}-1)}{\sin^2\left(\frac{1}{n}\right) + \ln\left(1+\frac{2}{n^2}\right)}.$
\end{enumerate}
2. \indent Si risolva il seguente integrale
$$\int \ln^2 x \, dx.$$
3. \indent Si disegni il grafico qualitativo della funzione
 \begin{equation*}
 f(x)=e^{-x} \sqrt{x-3}.
\end{equation*}
Si può trascurare lo studio della derivata seconda.\\ \\
4. \indent Si studi il carattere della serie
$$\sum_{n=1}^\infty \arctan \left(\dfrac{n+3}{n^2+4n}\right).$$
\textbf{Si risponda alle seguenti domande.} \\ \\
5. \indent  Data la funzione $f(x)=3x^2+4x$ si dica se soddisfa le ipotesi del teorema di Lagrange nell'intervallo $[0, 1]$. Si determini poi almeno un punto $x_0$ che soddisfa la tesi del teorema di Lagrange. \\ \\
6. \indent Si enunci e si dimostri la Formula Fondamentale del Calcolo Integrale. \\ \\
7. \indent Si presentino le definizioni di convergenza e assoluta convergenza di una serie, spiegando, anche attraverso esempi, i legami tra questi due concetti.\\ \\
%8. \indent Si dimostri che ogni funzione derivabile in un punto è continua in quel punto. Vale il viceversa? In caso contrario, si presenti un controesempio.



\end{document}
